\subsection{Summary}

This study offers a comprehensive examination of how Wi-Fi CSI can be leveraged for reliable human identification, demonstrating that carefully designed deep learning architectures—particularly convolutional and hybrid CNN–LSTM models—are essential for extracting discriminative gait signatures from complex radio-frequency measurements. The comparative evaluation conducted on 14 November established DenseNet1D as the strongest architecture among four candidates, owing to its dense connectivity, rich spatial feature extraction, and stable optimisation behaviour. The subsequent hyperparameter optimisation on 27 November further refined this architecture, resulting in a final model that achieved approximately 96\% accuracy on the held-out test set, supported by near-perfect macro precision, recall, and F1-scores. Together, these results reinforce the suitability of Wi-Fi CSI as a privacy-preserving biometric modality capable of high accuracy and robust generalisation across individuals.

While the current work demonstrates strong performance within controlled experimental conditions, it also highlights the importance of advancing the broader methodological ecosystem surrounding CSI-based sensing. Future progress will depend not only on improving model architectures but also on addressing key challenges in data diversity, cross-domain generalisation, and environmental variability \parencite{yang2023sensefi}. Developing adaptable models capable of maintaining accuracy across heterogeneous indoor settings, hardware platforms, and participant populations remains a pressing objective \parencite{li2025wifigeneralizability}. Equally important is the establishment of standardized protocols for CSI data collection, annotation, and benchmarking to support reproducible research and meaningful cross-study comparison \parencite{csibench2025}. Such standardisation would enable the creation of larger and more representative datasets, supporting the development of deep learning systems that maintain reliability in unconstrained real-world environments.

Several emerging research directions offer promising pathways for overcoming current limitations. Transfer learning and self-supervised feature pretraining could facilitate adaptation from synthetic CSI data to real-world deployment, alleviating the need for costly large-scale data collection \parencite{xu2023sslwifi}. Federated learning approaches may further support privacy-preserving model training across institutions by avoiding direct sharing of raw CSI data, thereby addressing both privacy and generalisation challenges \parencite{zhang2024fedwifi}. Additionally, integrating CSI with complementary passive sensing modalities—such as infrared, acoustic, or geomagnetic data—may yield multimodal fusion systems with improved robustness and situational awareness in complex environments \parencite{bandyopadhyay2025fedmagnet}.

Another important frontier lies in developing systems capable of operating reliably in multi-occupancy environments and varying architectural layouts, where overlapping human motion and dynamic multipath effects introduce significant ambiguity \parencite{rizk2025identifi}. Few-shot learning presents a promising solution to these challenges by enabling models to adapt quickly to new domains using limited labelled data, significantly reducing the need for exhaustive retraining \parencite{yang2023fewcross, li2025wififewshot}. The incorporation of generative adversarial techniques and meta-learning into few-shot frameworks has already shown substantial improvements in classification performance—sometimes exceeding 20\% in low-data conditions—by synthesizing realistic training samples and enhancing generalisation \parencite{bahadori2022rewis, yang2022metapose}. These advancements are increasingly vital as Wi-Fi sensing systems are pushed beyond controlled laboratory conditions into dynamic, diverse, and unpredictable environments.

Finally, this research contributes to a growing interdisciplinary movement that bridges wireless communication, machine learning, human–computer interaction, and biomechanics. Deeper investigation into the physiological underpinnings of gait—particularly the biomechanical signatures that manifest in CSI—may further strengthen the uniqueness and security of gait as a biometric identifier \parencite{mosharaf2024wifiid}. Such cross-disciplinary insights will be critical for developing robust, privacy-preserving, and practically deployable CSI-based identification systems capable of functioning reliably at scale.

In summary, while the present work demonstrates that deep learning models, especially DenseNet1D—can achieve high-accuracy gait identification using Wi-Fi CSI, the path toward widespread real-world adoption will require advances in data diversity, model interpretability, multi-domain generalisation, and ethical system design. Continued interdisciplinary research will be essential for realising the full potential of Wi-Fi CSI as a secure, unobtrusive, and pervasive sensing modality.


\begin{figure}[htbp]
\centering
\resizebox{\textwidth}{!}{%
\begin{tikzpicture}[
    node distance = 1.4cm and 1.8cm,
    >=Latex,
    font=\small,
    box/.style={
        draw,
        rounded corners,
        thick,
        align=center,
        minimum width=3.0cm,
        minimum height=0.9cm,
        fill=#1!15
    },
    bigbox/.style={
        draw,
        rounded corners,
        thick,
        align=center,
        minimum width=3.2cm,
        minimum height=1.0cm,
        fill=#1!10
    },
    dashedbox/.style={
        draw,
        rounded corners,
        thick,
        align=center,
        minimum width=3.2cm,
        minimum height=0.9cm,
        fill=gray!5,
        dashed
    },
    arrow/.style={->, thick}
]

%------------------------------------------------
% LEFT: SENSING ENVIRONMENT
%------------------------------------------------
\node[box=blue] (env) {Indoor Wi-Fi Setup};
\node[box=cyan, below=of env] (csi) {Raw CSI\\Amplitude + Phase};

%------------------------------------------------
% MID-LEFT: DATA PIPELINE
%------------------------------------------------
\node[box=green, right=of csi] (pre) {Preprocessing\\Parsing \& Windowing};

\node[box=yellow, above=of pre] (dataset) {CSI Dataset\\(30 Subjects, 12 Activities)};

\node[draw, rounded corners, thick, fit=(pre) (dataset), inner sep=5pt,
label=below:{\textbf{Data \& Processing}}] (datapipeline) {};

%------------------------------------------------
% MID: MLOps
%------------------------------------------------
\node[bigbox=orange, right=of dataset] (mlops) {MLOps / MLflow\\Run Tracking};

%------------------------------------------------
% RIGHT: MODELS
%------------------------------------------------
\node[bigbox=purple, below=of mlops] (models) {Models Compared:\\
CSILSTMNet\\DenseNet1D\\EfficientNet1D-LSTM\\MobileNetV3-1D-LSTM};

\node[dashedbox, right=of models] (best) {Best Model:\\\textbf{DenseNet1D}\\98.33\% Val Acc\\96\% Test Acc};

\node[box=red, below=of best] (results) {Evaluation\\Accuracy, F1, Recall};

\node[box=gray, below=of results] (apps) {Applications:\\Authentication\\Smart Homes\\Healthcare};

%------------------------------------------------
% MAIN FLOWS
%------------------------------------------------
\draw[arrow] (env) -- (csi);
\draw[arrow] (csi) -- (pre);
\draw[arrow] (dataset) -- (models);
\draw[arrow] (pre) -- (dataset);

% MLOps flows
\draw[arrow] (dataset) -- (mlops);
\draw[arrow] (models) -- (mlops);

% Right side
\draw[arrow] (models) -- (best);
\draw[arrow] (best) -- (results);
\draw[arrow] (results) -- (apps);

\end{tikzpicture}%
} % end resizebox
\caption{Graphical abstract of the CSI-based gait identification pipeline, illustrating sensing, data processing, model comparison, MLOps integration, and downstream applications.}
\label{fig:graphical_abstract_csi}
\end{figure}

